\documentclass[12pt]{beamer}
\usepackage[utf8]{inputenc}
\usepackage[T1]{fontenc}
\usepackage{lmodern}
\usetheme{Malmoe}
%\usetheme{Warsaw}
\usepackage{graphicx}
\usepackage{newverbs}

\newverbcommand{\tc}{\color{blue}}{}



\begin{document}
	\author{Lukáš Růžicka}
	\title{Modularita ve Fedoře 29}
	\subtitle{}
	\titlegraphic{\includegraphics[width=3cm]{logo.png}}
	\logo{\includegraphics[width=1cm]{logo.png}}
	\institute{Fedora 29 Release Party}
	\date{listopad 2018}
	\subject{Fedora 29}
	%\setbeamercovered{transparent}
	%\setbeamertemplate{navigation symbols}{}
	\begin{frame}[plain]
	\maketitle
\end{frame}

\section{Co je modularita?}

\begin{frame}
\frametitle{Co je vlastně modularita?}

\begin{itemize}
	\item nový plán na poskytování aplikací
	\item řeší problémy tradiční distribuce software
	\item větší přizpůsobení se uživatelským potřebám
\end{itemize}
\end{frame}

\begin{frame}
\frametitle{Co musíte vědět o Modularitě?}

\begin{itemize}
	\item poskytuje větší flexibilitu v distribuci aplikací
	\item umožňuje řešit nové uživatelské případy
	\item má vždy přednost před běžnými repozitáři
	\item neumožňuje instalaci více verzí jedné aplikace na jednom stroji
	\item neřeší (pokud vím) závislosti dalších aplikací
	\item jestliže nejsou dostupná modulární repa, systém o modularitě neví a při operacích s balíčky může dojít k~nečekaným překvapením 
\end{itemize}
\end{frame}


\begin{frame}
\frametitle{Modularita pro všechny}
Do \textbf{Fedory 29} přichází globální podpora pro modularitu, kterou jsme zatím mohli potkat pouze ve \textbf{Fedora Server 28}. 

Nyní ve verzi 29, modularita:

\vspace{5pt}

\begin{itemize}
	\item je zapnutá napříč spiny
	\item není pro použití Fedory klíčová
	\item pokrývá pouze malé spektrum aplikací (cca 30)
	\item je rámcově funkční
\end{itemize}	
\end{frame}

\begin{frame}
\frametitle{Co pro vás může modularita udělat?}
Nejdiskutovanější případy použití, které modularita pro uživatele řeší, jsou:

\vspace{5pt}

\begin{itemize}
	\item standardní vývoj Fedory je příliš \textit{pomalý} nebo \textit{rychlý}
	\item lepší použití v kontejnerovém prostředí
	\item snadnější správa balíčků
\end{itemize}
\end{frame}

\section{Správa software a modularita}

\begin{frame}
\frametitle{Tradiční pojetí balíčků}
V tradičním pojetí se balíčky vytvářejí z jedné verze zdrojových kódů pro určitou verzi Fedory. V pozdějších verzích Fedory se objevují novější balíčky.
\begin{center}
	\includegraphics[width=10cm]{traditional}
\end{center}
\end{frame}

\begin{frame}
\frametitle{Modulární pojetí balíčků}
V \textit{modulárním} pojetí se balíčky vytvářejí z jedné verze zdrojových kódů pro různé verze Fedory. Jednotlivé verze balíčků mohou být dostupné v několika verzích Fedory.
\begin{center}
	\includegraphics[width=10cm]{modular}
\end{center}
\end{frame}

\begin{frame}
\frametitle{Přehled modulárního OS}

\begin{center}
	\includegraphics[width=10cm]{overview}
\end{center}
\end{frame}

\begin{frame}
\frametitle{Jedna verze balíčků, více verzí Fedory}
Modul může být nastaven tak, aby se balíček přeložil pro několik verzí Fedory najednou. 

\begin{center}
	\includegraphics[width=7cm]{forversion}
\end{center}
\end{frame}

\begin{frame}
\frametitle{Více verzí balíčků, jedna verze Fedory}
Moduly lze využít i tak, že pro jednu verzi Fedory se přeloží několik různých verzí příslušného software.

\begin{center}
	\includegraphics[width=7cm, height=3.5cm]{multiversion}
\end{center}

Modulární přístup však nedovoluje mít v systému více verzí jedné aplikace, bez použití dalších technologií (kontejner).

\end{frame}

\begin{frame}
\frametitle{Složitější strategie}
Moduly lze také přeložit oproti dalším modulům a tak je možné vyřešit i daleko komplexnější požadavky uživatelů.

\begin{center}
	\includegraphics[width=8cm]{combined}
\end{center}
\end{frame}

\begin{frame}
\frametitle{Organizace repozitářů}

\begin{enumerate}
	\item tradiční repozitář
	\item modulární repozitář (AppStream)
\end{enumerate}

\begin{center}
	\includegraphics[width=10cm]{repos}
\end{center}
\end{frame}

\begin{frame}
\frametitle{Modulární repozitáře ve Fedoře 29}

\begin{enumerate}
	\item fedora-modular.repo
	\item fedora-updates-modular.repo
	\item fedora-updates-testing-modular.repo
\end{enumerate}

\begin{center}
	\includegraphics[width=9cm]{repos}
\end{center}
\end{frame}

\begin{frame}[fragile]
\frametitle{Modulární repozitáře ve Fedoře 29}

\begin{enumerate}
	\item standardně zapnuté pro všechny
	\item není třeba nic víc řešit
 	\item umístění v \tc|/etc/yum.repos.d/|, pokud někdo chce měnit jejich nastavení
\end{enumerate}
\end{frame}

\section{Používání modulárního obsahu}

\begin{frame}
\frametitle{Práce s modulem}

K používání modulárního obsahu existuje několik operací:

\begin{enumerate}
	\item zobrazit seznam modulů
	\item povolit modul
	\item nainstalovat modul
	\begin{itemize}
		\item stream
		\item profile
	\end{itemize}
	\item změnit \textit{stream} modulu
	\item odinstalovat modul
	\item zakázat modul
\end{enumerate}
\end{frame}

\begin{frame}[fragile]
\frametitle{Zobrazení seznamu dostupných modulů}

Zobrazení seznamu všech dostupných modulů:

\begin{center}
	\tc|dnf module list|
\end{center}

Informace o konkrétním modulu:

\begin{center}
	\tc|dnf module list <modulename>|
\end{center}
\end{frame}

\begin{frame}[fragile]
\frametitle{Další možnosti zobrazení}

Moduly si lze také zobrazit podle jejich stavu. Každý modul může být buď:

\begin{itemize}
	\item zakázaný (\tc|--disabled|)
	\item povolený (\tc|--enabled|)
	\item nainstalovaný (\tc|--installed|)
\end{itemize}

Přepínače uvedené v závorkách jsou přepínače příkazu \tc|dnf module list|.
\end{frame}

\begin{frame}[fragile]
\frametitle{Informace o modulu}
Ve výpisu je pro každý modul uvedeno:

\begin{verbatim}
mariadb  10.1  client, server, default  #info
\end{verbatim}

\begin{itemize}
	\item název -- většinou jméno projektu (\tc|mariadb|)
	\item stream -- verze projektu (\tc|10.1|)
	\item profil -- sada balíčků (\tc|client, server, default|)
	\item info -- krátký popis modulu
\end{itemize}
\end{frame}

\begin{frame}[fragile]
\frametitle{Povolení modulu}
Povolení modulu umožní správci balíčků, aby s tímto modulem pracoval. 

Následující příkaz povolí modul \tc|perl| ve verzi \tc|5.26|:

\begin{center}
	\tc|dnf module enable perl:5.26|
\end{center}

Takto povolený modul lze pak instalovat jako běžný balíček, tedy \tc|dnf install perl|.
\end{frame}

\begin{frame}[fragile]
\frametitle{Zakázání modulu}
Zakázaný modul bude správce balíčků ignorovat. 

Následující příkaz zakáže modul \tc|perl| ve verzi \tc|5.26|:

\begin{center}
	\tc|dnf module disable perl:5.26|
\end{center}

\tc|dnf install perl| tak nainstaluje verzi jazyka \tc|perl|, která se nachází ve standardním repozitáři, pokud nebude povolen jiný \textit{stream} téhož modulu.
\end{frame}

\begin{frame}[fragile]
\frametitle{Instalace modulu}

\begin{itemize}
	\item Modul lze instalovat i explicitním příkazem.
	\item Není potřeba jej předem povolovat, povolí se během instalace.
	\item Takto je možné zvolit si nejen konkrétní \textit{stream}, ale také \textit{profil} modulu.
\end{itemize}

\begin{center}
	\tc|dnf module install mariadb:10.1/client|
\end{center}

\end{frame}

\begin{frame}[fragile]
\frametitle{Přednastavené hodnoty (default stream a profil)}

V budoucnu bude mít každý modul definován svůj základní \textit{stream} a \textit{profil}. Práce s moduly by měla být ještě pohodlnější. Funkčnost by měla odpovídat následujícímu:

\begin{itemize}
	\item \tc|dnf install nodejs| nainstaluje RPM balíček \tc|nodejs| ze základního streamu.
	\item \tc|dnf module install nodejs| nainstaluje základní profil a stream.
	\item \tc|dnf module install nodejs:9| nainstaluje základní profil ze streamu \tc|9|.
\end{itemize}
\end{frame}

\begin{frame}[fragile]
\frametitle{Jak vypnout Modularitu pomocí konfiguračních souborů?}

\begin{enumerate}
	\item Najděte konfigurační soubory modulárních repozitářů.
	
	\tc|ls /etc/yum.repos.d/*modular*|
	\item Otevřete jeden z nalezených souborů.
	
	\tc|vi /etc/yum.repos.d/fedora-modular.repo|
	\item V základní sekci (zde \tc|[fedora-modular]|) najděte položku \tc|enabled=1|.
	\item Změňte její hodnotu na \tc|enabled=0|.
	\item Zopakujte postup pro všechny repozitáře, které chcete zakázat.
\end{enumerate}
\end{frame}

\begin{frame}[fragile]
\frametitle{Jak vypnout Modularitu pomocí DNF?}

\begin{enumerate}
	\item Zobrazte si dostupné povolené repozitáře.
	
	\tc|dnf repolist --enabled|
	\item Ze seznamu si vyberte nějaký modulární repozitář (v~názvu je slovo \textit{modular}), např. \tc|fedora-modular|.
	\item Zakažte tento repozitář.
	
	\tc|dnf config-manager --set-disabled fedora-modular|
	\item Zopakujte postup pro všechny repozitáře, které chcete zakázat.
\end{enumerate}
\end{frame}

\begin{frame}[fragile]
\frametitle{Více informací}
	Dokumentace o současném stavu se nachází na \tc|https://docs.fedoraproject.org/en-US/modularity/|.
	
	Odsud jsem si vypůjčil diagramy.
\end{frame}

\begin{frame}
\frametitle{Otázky a odpovědi}

Ptejte se, protože

\begin{center}
	\Huge{Líná huba, holé neštěstí.}
\end{center}
\end{frame}

\begin{frame}
\frametitle{Díky za pozornost}
\begin{center}
	\Huge{a užijte si Fedoru 29!}
\end{center}
\end{frame}



\end{document}
